\chapter{Thực nghiệm và đánh giá mô hình}

\section{Môi trường và dữ liệu đánh giá}

Các thực nghiệm được thực hiện trên cùng một môi trường cục bộ của dự án để đảm bảo dữ liệu, cấu hình huấn luyện và kết quả đánh giá có thể truy vết lại. Chương này sử dụng các artifact cuối cùng trong thư mục \texttt{results/}, đặc biệt là file tổng hợp \texttt{results/final\_figures/final\_metrics\_for\_report.csv}.

\begin{table}[H]
  \centering
  \caption{Môi trường phần mềm sử dụng trong thực nghiệm}
  \begin{tabular}{p{4.2cm} p{8.8cm}}
    \toprule
    \textbf{Thành phần} & \textbf{Giá trị} \\
    \midrule
    Ngôn ngữ lập trình & \texttt{Python 3.10.0}. \\
    Thư viện xử lý dữ liệu & \texttt{pandas 2.3.3}, \texttt{numpy 2.2.6}. \\
    Thư viện học máy & \texttt{scikit-learn 1.6.1}, \texttt{xgboost 3.2.0}. \\
    Thư viện trực quan hóa & \texttt{matplotlib 3.10.8}. \\
    Dữ liệu xử lý & \texttt{data/processed/raw/}, \texttt{data/processed/scaled/} và \texttt{data/processed/stats/}. \\
    Thư mục kết quả & \texttt{results/1h/}, \texttt{results/24h/} và \texttt{results/final\_figures/}. \\
    \bottomrule
  \end{tabular}
  \label{tab:experiment_environment}
\end{table}

Bảng \ref{tab:experiment_environment} cho thấy các thực nghiệm được thực hiện bằng bộ thư viện thống nhất. Báo cáo không sử dụng thông tin phần cứng để giải thích kết quả vì cấu hình CPU, RAM hoặc GPU không được lưu trong artifact của lần chạy.

\section{Thiết lập thực nghiệm cuối cùng}

Thực nghiệm được thiết kế cho hai horizon: \texttt{1h} và \texttt{24h}. Cả hai horizon đều dùng tập test năm 2025, tức dữ liệu không tham gia vào quá trình học tham số. Mô hình ANFIS cuối cùng sử dụng Core Features gồm \texttt{apparent\_temperature}, \texttt{humidity}, \texttt{hour\_sin}, \texttt{hour\_cos}, \texttt{occupancy\_level}, \texttt{load\_lag\_1} và \texttt{load\_lag\_24}.

\begin{table}[H]
  \centering
  \caption{Cấu hình ANFIS dùng cho kết quả cuối cùng}
  \begin{tabular}{p{2.1cm} p{4.4cm} p{2.7cm} r r r}
    \toprule
    \textbf{Horizon} & \textbf{Run ID} & \textbf{Target} & \textbf{Inputs} & \textbf{\texttt{n\_mfs}} & \textbf{Luật} \\
    \midrule
    \texttt{1h} & \texttt{1h\_20260617\_151950} & \texttt{target\_1h} & 7 & 2 & 128 \\
    \texttt{24h} & \texttt{24h\_20260617\_151950} & \texttt{target\_24h} & 7 & 2 & 128 \\
    \bottomrule
  \end{tabular}
  \label{tab:final_anfis_config}
\end{table}

Bảng \ref{tab:final_anfis_config} thể hiện cấu hình cuối của ANFIS sau khi bổ sung \texttt{load\_lag\_1}. Với 7 đặc trưng đầu vào và 2 hàm thuộc Gaussian cho mỗi đặc trưng, số luật mờ là $2^7=128$, giúp mô hình giữ được độ phức tạp ở mức có thể kiểm soát.

\begin{table}[H]
  \centering
  \caption{Thiết lập dữ liệu theo từng horizon}
  \begin{tabular}{p{3cm} r r r p{3cm}}
    \toprule
    \textbf{Horizon} & \textbf{Train-fit} & \textbf{Validation} & \textbf{Test} & \textbf{Khoảng test} \\
    \midrule
    \texttt{1h} & 105\,024 & 35\,136 & 34\,944 & Năm 2025 \\
    \texttt{24h} & 105\,024 & 35\,136 & 34\,944 & Năm 2025 \\
    \bottomrule
  \end{tabular}
  \label{tab:final_data_split}
\end{table}

Bảng \ref{tab:final_data_split} cho thấy hai horizon dùng cùng cách chia dữ liệu theo thời gian. Việc tách train, validation và test theo mốc thời gian giúp hạn chế rò rỉ dữ liệu từ tương lai vào quá trình huấn luyện.

\section{Các mô hình so sánh}

ANFIS là mô hình chính của đề tài và sử dụng Core Dataset. Các mô hình \texttt{Linear Regression}, \texttt{Decision Tree}, \texttt{Random Forest} và \texttt{XGBoost} được huấn luyện lại trên Extended Dataset để làm baseline so sánh.

\begin{table}[H]
  \centering
  \caption{Các mô hình sử dụng trong thực nghiệm}
  \begin{tabularx}{\textwidth}{p{3.6cm} p{2.8cm} X}
    \toprule
    \textbf{Mô hình} & \textbf{Dữ liệu} & \textbf{Vai trò trong thực nghiệm} \\
    \midrule
    \texttt{ANFIS} & \texttt{core} & Mô hình chính của đề tài, dùng 7 Core Features, hàm thuộc Gaussian và hệ quả Sugeno bậc nhất. \\
    \texttt{Linear Regression} & \texttt{extended} & Baseline tuyến tính, dùng để kiểm tra mức hiệu quả tối thiểu khi quan hệ được giả định tuyến tính. \\
    \texttt{Decision Tree} & \texttt{extended} & Baseline phi tuyến đơn giản, có khả năng học các ngưỡng quyết định trong dữ liệu. \\
    \texttt{Random Forest} & \texttt{extended} & Mô hình ensemble nhiều cây, thường ổn định hơn cây đơn trên dữ liệu bảng. \\
    \texttt{XGBoost} & \texttt{extended} & Mô hình boosting mạnh, khai thác tốt các tương tác phi tuyến giữa đặc trưng mở rộng. \\
    \bottomrule
  \end{tabularx}
  \label{tab:models_used_final}
\end{table}

Bảng \ref{tab:models_used_final} cho thấy sự khác nhau giữa mô hình chính và baseline: ANFIS ưu tiên tập đặc trưng nhỏ để giữ khả năng diễn giải, trong khi các baseline dùng tập Extended để tối ưu hiệu năng dự báo.

\section{Kết quả dự báo 1 giờ tới}

\begin{table}[H]
  \centering
  \caption{Kết quả cuối cùng cho horizon \texttt{1h} trên tập test}
  \resizebox{\textwidth}{!}{%
  \begin{tabular}{p{4.2cm} r r r r}
    \toprule
    \textbf{Model} & \textbf{MAE} & \textbf{RMSE} & \textbf{MAPE (\%)} & \textbf{$R^2$} \\
    \midrule
    \texttt{Linear Regression} & 0.205078 & 0.301607 & 33.430 & 0.735643 \\
    \texttt{Decision Tree} & 0.106042 & 0.180408 & 17.262 & 0.905415 \\
    \texttt{Random Forest} & \textbf{0.084448} & \textbf{0.153935} & \textbf{13.905} & \textbf{0.931137} \\
    \texttt{XGBoost} & 0.086653 & 0.154348 & 14.567 & 0.930767 \\
    \texttt{ANFIS} & 0.110881 & 0.187050 & 18.140 & 0.898323 \\
    \bottomrule
  \end{tabular}
  }
  \label{tab:results_1h}
\end{table}

Bảng \ref{tab:results_1h} cho thấy ANFIS đạt RMSE 0.187050 kWh và $R^2$ 0.898323 ở mốc 1 giờ tới. Kết quả này tốt hơn rõ rệt so với \texttt{Linear Regression}, nhưng \texttt{Random Forest} và \texttt{XGBoost} vẫn cho sai số thấp hơn nhờ khai thác Extended Dataset.

\begin{figure}[H]
  \centering
  \includegraphics[width=\textwidth]{../results/1h/plots/1h_20260617_151950_actual_vs_predicted.png}
  \caption{Actual vs Predicted của ANFIS cho horizon \texttt{1h}}
  \label{fig:anfis_actual_pred_1h}
\end{figure}

Hình \ref{fig:anfis_actual_pred_1h} cho thấy dự báo ANFIS nhìn chung bám theo xu hướng phụ tải thực tế trong giai đoạn minh họa. Một số vùng có spike hoặc biến động nhanh vẫn bị làm mượt, đây là nguyên nhân khiến RMSE cao hơn các mô hình ensemble.

\begin{figure}[H]
  \centering
  \includegraphics[width=0.82\textwidth]{../results/1h/plots/1h_20260617_151950_residuals.png}
  \caption{Phân phối residual của ANFIS cho horizon \texttt{1h}}
  \label{fig:anfis_residual_1h}
\end{figure}

Hình \ref{fig:anfis_residual_1h} cho thấy residual tập trung quanh 0, phản ánh mô hình không bị lệch hệ thống quá lớn trên tập test. Phần đuôi của phân phối vẫn tồn tại do các thời điểm phụ tải tăng hoặc giảm đột ngột trong dữ liệu mô phỏng.

\section{Kết quả dự báo 24 giờ tới}

\begin{table}[H]
  \centering
  \caption{Kết quả cuối cùng cho horizon \texttt{24h} trên tập test}
  \resizebox{\textwidth}{!}{%
  \begin{tabular}{p{4.2cm} r r r r}
    \toprule
    \textbf{Model} & \textbf{MAE} & \textbf{RMSE} & \textbf{MAPE (\%)} & \textbf{$R^2$} \\
    \midrule
    \texttt{Linear Regression} & 0.172051 & 0.250550 & 27.501 & 0.817623 \\
    \texttt{Decision Tree} & 0.148690 & 0.230826 & 24.107 & 0.845207 \\
    \texttt{Random Forest} & 0.134404 & 0.213531 & 21.496 & 0.867535 \\
    \texttt{XGBoost} & \textbf{0.133447} & \textbf{0.210512} & \textbf{21.422} & \textbf{0.871254} \\
    \texttt{ANFIS} & 0.145392 & 0.231865 & 23.355 & 0.843810 \\
    \bottomrule
  \end{tabular}
  }
  \label{tab:results_24h}
\end{table}

Bảng \ref{tab:results_24h} cho thấy với horizon \texttt{24h}, ANFIS đạt RMSE 0.231865 kWh và $R^2$ 0.843810. Sai số cao hơn horizon \texttt{1h}, phù hợp với bản chất khó hơn của bài toán dự báo xa hơn; trong nhóm baseline, \texttt{XGBoost} đạt kết quả tốt nhất.

\begin{figure}[H]
  \centering
  \includegraphics[width=\textwidth]{../results/24h/plots/24h_20260617_151950_actual_vs_predicted.png}
  \caption{Actual vs Predicted của ANFIS cho horizon \texttt{24h}}
  \label{fig:anfis_actual_pred_24h}
\end{figure}

Hình \ref{fig:anfis_actual_pred_24h} cho thấy ANFIS vẫn nắm bắt được chu kỳ ngày của phụ tải, nhưng sai lệch tại các vùng cao điểm rõ hơn so với horizon \texttt{1h}. Điều này phản ánh mức bất định tăng khi mô hình phải dự báo trước một ngày.

\begin{figure}[H]
  \centering
  \includegraphics[width=0.82\textwidth]{../results/24h/plots/24h_20260617_151950_residuals.png}
  \caption{Phân phối residual của ANFIS cho horizon \texttt{24h}}
  \label{fig:anfis_residual_24h}
\end{figure}

Hình \ref{fig:anfis_residual_24h} cho thấy residual vẫn tập trung quanh 0 nhưng phân tán rộng hơn so với horizon \texttt{1h}. Điều này phù hợp với kết quả RMSE lớn hơn và cho thấy bài toán dự báo 24 giờ tới nhạy hơn với các biến động hành vi trong dữ liệu.

\section{So sánh mô hình bằng RMSE và $R^2$}

\begin{figure}[H]
  \centering
  \begin{subfigure}{0.48\textwidth}
    \centering
    \includegraphics[width=\textwidth]{../results/final_figures/rmse_comparison_1h_final.png}
    \caption{Horizon \texttt{1h}}
  \end{subfigure}
  \hfill
  \begin{subfigure}{0.48\textwidth}
    \centering
    \includegraphics[width=\textwidth]{../results/final_figures/rmse_comparison_24h_final.png}
    \caption{Horizon \texttt{24h}}
  \end{subfigure}
  \caption{So sánh RMSE giữa các mô hình trong thực nghiệm cuối cùng}
  \label{fig:rmse_comparison_final}
\end{figure}

Hình \ref{fig:rmse_comparison_final} cho thấy ANFIS không phải mô hình có RMSE thấp nhất. Ở horizon \texttt{1h}, \texttt{Random Forest} tốt nhất; ở horizon \texttt{24h}, \texttt{XGBoost} tốt nhất. Tuy vậy, ANFIS vẫn đạt sai số cạnh tranh khi chỉ sử dụng Core Features.

\begin{figure}[H]
  \centering
  \begin{subfigure}{0.48\textwidth}
    \centering
    \includegraphics[width=\textwidth]{../results/final_figures/r2_comparison_1h_final.png}
    \caption{Horizon \texttt{1h}}
  \end{subfigure}
  \hfill
  \begin{subfigure}{0.48\textwidth}
    \centering
    \includegraphics[width=\textwidth]{../results/final_figures/r2_comparison_24h_final.png}
    \caption{Horizon \texttt{24h}}
  \end{subfigure}
  \caption{So sánh $R^2$ giữa các mô hình trong thực nghiệm cuối cùng}
  \label{fig:r2_comparison_final}
\end{figure}

Hình \ref{fig:r2_comparison_final} cho thấy các mô hình ensemble đạt $R^2$ cao nhất ở cả hai horizon. ANFIS vẫn giải thích được phần lớn biến thiên của phụ tải, đặc biệt ở horizon \texttt{1h}, nhưng còn khoảng cách so với các baseline mạnh trên Extended Dataset.

\section{Các cải tiến so với bản nộp lần 1}

So với bản nộp lần 1, nhóm đã hoàn thiện cả dữ liệu đầu vào, thực nghiệm và phần phân tích kết quả:

\begin{itemize}[leftmargin=1.2cm]
  \item Bổ sung \texttt{load\_lag\_1} vào Core Features để mô hình khai thác thêm thông tin phụ tải của một giờ gần nhất.
  \item Huấn luyện lại ANFIS cho cả hai horizon và train lại các mô hình baseline để có bảng so sánh cuối cùng.
  \item Cập nhật bảng metric cuối cùng cho \texttt{1h} và \texttt{24h} từ file CSV tổng hợp trong \texttt{results/final\_figures/}.
  \item Bổ sung biểu đồ so sánh mô hình theo RMSE và $R^2$.
  \item Bổ sung phân tích residual để đánh giá phân phối sai số của ANFIS.
  \item Bổ sung nhận xét về hạn chế của dữ liệu phụ tải mô phỏng.
\end{itemize}

\begin{table}[H]
  \centering
  \caption{So sánh kết quả ANFIS trước và sau cải tiến}
  \resizebox{\textwidth}{!}{%
  \begin{tabular}{p{2cm} r r r r r r}
    \toprule
    \textbf{Horizon} & \textbf{RMSE lần 1} & \textbf{RMSE cuối} & \textbf{$\Delta$ RMSE} & \textbf{$R^2$ lần 1} & \textbf{$R^2$ cuối} & \textbf{$\Delta R^2$} \\
    \midrule
    \texttt{1h} & 0.1967 & 0.187050 & -0.009650 & 0.8876 & 0.898323 & +0.010723 \\
    \texttt{24h} & 0.2357 & 0.231865 & -0.003835 & 0.8385 & 0.843810 & +0.005310 \\
    \bottomrule
  \end{tabular}
  }
  \label{tab:anfis_improvement}
\end{table}

Bảng \ref{tab:anfis_improvement} cho thấy ANFIS cải thiện ở cả hai horizon: RMSE giảm và $R^2$ tăng so với bản nộp lần 1. Mức cải thiện rõ hơn ở horizon \texttt{1h}, phù hợp với vai trò của \texttt{load\_lag\_1} trong việc cung cấp thông tin phụ tải rất gần thời điểm dự báo.

\section{Nhận xét chung}

Kết quả tổng thể cho thấy mốc dự báo \texttt{1h} dễ hơn mốc \texttt{24h}. Với ANFIS, RMSE tăng từ 0.187050 kWh ở horizon \texttt{1h} lên 0.231865 kWh ở horizon \texttt{24h}; $R^2$ giảm từ 0.898323 xuống 0.843810. Xu hướng này phản ánh độ bất định tăng khi khoảng cách dự báo xa hơn.

ANFIS đạt kết quả cạnh tranh và có ưu điểm quan trọng là khả năng diễn giải thông qua membership functions và luật mờ. Tuy nhiên, nếu tiêu chí chính là RMSE thấp nhất, \texttt{Random Forest} và \texttt{XGBoost} đang hiệu quả hơn trong thực nghiệm hiện tại vì được huấn luyện trên Extended Dataset với nhiều đặc trưng mở rộng hơn.

Cần lưu ý rằng dữ liệu phụ tải trong đề tài là dữ liệu mô phỏng từ profile hộ gia đình, chưa phải dữ liệu công tơ điện thực tế. Vì vậy, kết quả nên được hiểu là đánh giá quy trình và hướng tiếp cận trong phạm vi học phần, chưa phải kết luận vận hành trên dữ liệu thực tế.
